\chapter{Introduction}
Lung disease remains a leading cause of global morbidity and mortality \cite{murray2022global}. Research shows that early
detection of pulmonary diseases significantly improves patient outcomes. For this, radiological imaging tools
such as Chest X-Ray (CXR) and Computed Tomography (CT), plays an important role. However, the increasing volume of
such imaging data increases radiologists' workload, leading to delays and potential errors in diagnosis.
This has motivated the development of automated Computer-Aided Diagnosis (CAD) systems, which use computational algorithms
and Artificial Intelligence (AI) to assist in tasks such as abnormality detection, classification, and segmentation \cite{doi2007computer}.

\section{Background}

Developing a unified pipeline for multi-class abnormality classification and grounded localization
addresses two of the most critical imperatives in clinical AI adoption: diagnostic accuracy and explainability.
Specifically, utilizing dense pixel-level segmentation as the mechanism for localization forces the architecture to
provide voxel-level visual evidence for its diagnostic claims. This ensures that the model is accountable and its
claims are supported by clinical evidence.

Our research mainly focuses on the development of a CAD system using 3 Dimensional (3D) CT data. While CXR is commonly
used as the first imaging test due to its low cost and low radiation exposure, studies show that their sensitivity
to abnormailities is limited \cite{borakati2020diagnostic}. Comparatively, CT scans produce high resolution
cross-sectional slices that helps detect smaller, early-stage, and subtle lung pathologies that are
frequently missed by standard CXR.

However, CT analysis is complex and time-consuming, requiring radiologists to review hundreds of slices per scan,
which are often inconsistent due to artifacts, anatomical variations and different scanners/protocols.
Furthermore, different anatomical structures can share identical grayscale values. Telling them apart
requires expert clinical judgment which is often subjective. High sensitivity can lead to detection of 
clinically insignificant abnormalities, while visual fatigue can lead to missed findings \cite{mccreadie2009eight}.

To address these challenges, recent research and clinical deployments of AI-based
CAD systems are fundamentally changing how radiologists approach CT analysis. AI models such as 3D
Convolutional Neural Networks (CNNs) and Vision Transformers (ViTs) have shown to reduce interpretation time, and improve
lesion detection sensitivity compared to traditional manual methods \cite{LITJENS201760}.
Many of these high performing models are Deep Learning based, which are criticized for their black box nature,
where their internal reasoning is difficult to interpret \cite{tjoa2021survey}. This lack of interpretability is a significant barrier for
clinical adoption, as radiologists require evidence to support a model's diagnostic claims.

This has motivated the development of explainable AI (XAI) methods such as Gradient-weighted Class Activation Mapping (Grad-CAM)
that helps to visualize exactly which parts of an image a model focuses on when making a specific prediction.
Such XAI methods are highly valuable for building clinical trust in deep learning models. However, they
often lack the fine-grained, pixel-level precision required to identify tiny lesions,
sometimes highlighting incorrect or confounding features \cite{ennab2025advancing}.

Recent advances in sparse representation models \cite{liyanage2022covid, wesp2026sparse} such as Sparse Autoencoders
and Dictionary Learning have shown to provide
inherent interpretability in medical imaging. This stems from their ability to reconstruct complex, high-dimensional
data as a linear combination of "atoms" or features. This allows the use of efficient algorithms to easily map the
output sparse domain back to the original input feature space, providing a linear relationship between the model's
decision and anatomical structures.
Our work builds on these advances by applying dictionary learning to 3D Chest CT data for multi-class abnormality classification
and grounded localization.

\section{Problem Statement}
The clinical utility of Deep Learning based approaches is hindered by a lack of interpretability. While tools like Grad-CAM
provide post-hoc explanations to model decisions, these are often insufficient for medical professionals who require a granular
understanding of the model's decision making process. In high stakes environments, a model that produces a diagnosis without
transparency in its reasoning is difficult to integrate into clinical workflows, as it lacks the accountability necessary for
legal and ethical compliance. \cite{tjoa2021survey}

\section{Objectives and Scope}
The primary goal of this study is to address the above mentioned limitations of existing AI-based approaches for abnormality
classification and localization in chest CTs, by developing an inherently interpretable dictionary learning framework.
This approach aims to enable transparent, anatomically grounded decision making, where each classification outcome
can be traced back to specific pathological patterns. This aligns with clinical reasoning and facilitates validation and
adoption in radiological workflows. The specific objectives are as follows:

\subsection{Objectives}
\begin{enumerate}
	\item Implement a Dictionary Learning based CAD system using a publicly available annotated 3D chest CT dataset.
	\item Map dictionary atoms back to anatomical features and visualize their spatial contribution maps.
	\item Develop a unified evaluation framework for assessing the performance and interpretability of the proposed approach.
	\item Benchmark model performance against SotA classification and segmentation models on the same dataset.
\end{enumerate}

\subsection{Scope}
The scope of this study focuses on the development and evaluation of a dictionary learning based approach for classification
and grounded localization of lung abnormalities using volumetric 3D CT data, with the following limitations:

\begin{itemize}
	\item This study is limited to 2 thoracic abnormalities: Pulmonary Nodules, and Ground-Glass Opacities, based on expert radiologist recommendation and dataset availbility.
	\item This study will utilize a single publicly available dataset for training and evaluation, which may limit generalizability to other populations and scanner types.
	\item The proposed approach will be limited only to interpretable methods, and will not explore approaches that provide better performance for less explainability.
\end{itemize}


%%%%%%%%%%%%%%%%%%%%%%%%%%%%%%%%%%%%%%%%%%%%%%%%%%%%%%%%%%%%%%%%%%%%%%%%%%%%%%%%%%%%%%%%%%%%%%%%%%
%%%%%%%%%%% END OF FILE
%%%%%%%%%%%%%%%%%%%%%%%%%%%%%%%%%%%%%%%%%%%%%%%%%%%%%%%%%%%%%%%%%%%%%%%%%%%%%%%%%%%%%%%%%%%%%%%%%%
